\documentclass[../../main.tex]{subfiles}% 注意这里的文件路径不能用 ./main.tex ,否则用latexmk编译子文件会报错
\graphicspath{{\subfix{./image/}}{\subfix{../../image/}}} % 指定图片引用路径,后续可以直接使用图片文件名
% 如果图片存放在上级目录,则使用 ../../image/ 作为备用路径

% 例如:
% \begin{figure}[H]
% \centering
% \includegraphics[scale=0.3]{图.png}
% \caption{}
% \label{figure:图}
% \end{figure}
% 注意:上述\label{}一定要放在\caption{}之后,否则引用图片序号会只会显示??.

\begin{document}

\section{主方向和主曲率的计算}

由法曲率的Euler公式，主曲率是法曲率的最大值与最小值，主方向与主曲率对应Weingarten映射的特征方向与特征值。

设正则参数曲面$S:\boldsymbol{r}=\boldsymbol{r}(u,v)$，$\delta\boldsymbol{r}=\boldsymbol{r}_u\delta u+\boldsymbol{r}_v\delta v$为主方向，即$(\delta u,\delta v)\neq(0,0)$，存在实数$\lambda$使得
\begin{align}
W(\delta\boldsymbol{r})=\lambda\delta\boldsymbol{r}.\label{eq:weingarten_eigen_eq}
\end{align}
代入Weingarten映射定义$W(\delta\boldsymbol{r})=-(\boldsymbol{n}_u\delta u+\boldsymbol{n}_v\delta v)$，得
\begin{align}
-(\boldsymbol{n}_u\delta u+\boldsymbol{n}_v\delta v)=\lambda(\boldsymbol{r}_u\delta u+\boldsymbol{r}_v\delta v).\label{eq:weingarten_eigen_expand}
\end{align}
将上式分别与$\boldsymbol{r}_u,\boldsymbol{r}_v$作内积，得线性方程组
\begin{align}
\begin{cases}
(L-\lambda E)\delta u+(M-\lambda F)\delta v=0,\\
(M-\lambda F)\delta u+(N-\lambda G)\delta v=0.
\end{cases}\label{eq:principal_dir_linear_eq}
\end{align}
该齐次方程组存在非零解的充要条件为系数行列式为零，即
\begin{align}
\begin{vmatrix}
L-\lambda E & M-\lambda F\\
M-\lambda F & N-\lambda G
\end{vmatrix}=0.\label{eq:principal_curvature_char_eq}
\end{align}
展开行列式得到$\lambda$的二次方程
\begin{align}
\lambda^2(EG-F^2)-\lambda(LG-2MF+NE)+(LN-M^2)=0,\label{eq:principal_curvature_quad_eq}
\end{align}
整理为
\begin{align}
\lambda^2-\frac{LG-2MF+NE}{EG-F^2}\lambda+\frac{LN-M^2}{EG-F^2}=0.\label{eq:principal_curvature_quad_simp}
\end{align}
其判别式
\begin{align*}
(LG-2MF+NE)^2-4(EG-F^2)(LN-M^2)\geqslant0,
\end{align*}
故方程必有两个实根$\kappa_1,\kappa_2$，即曲面的两个主曲率。

\begin{definition}
曲面的\textbf{平均曲率}定义为$H=\frac{1}{2}(\kappa_1+\kappa_2)$，\textbf{Gauss曲率（总曲率）}定义为$K=\kappa_1\kappa_2$。
\end{definition}

由二次方程根与系数的关系，可得
\begin{align}
\kappa_1+\kappa_2=2H=\frac{LG-2MF+NE}{EG-F^2},\label{eq:mean_curvature_formula}\\
\kappa_1\kappa_2=K=\frac{LN-M^2}{EG-F^2}.\label{eq:gauss_curvature_formula}
\end{align}
主曲率满足方程
\begin{align}
\lambda^2-2H\lambda+K=0,\label{eq:principal_curvature_hk_eq}
\end{align}
解得
\begin{align}
\kappa_1=H+\sqrt{H^2-K},\quad \kappa_2=H-\sqrt{H^2-K}.\label{eq:principal_curvature_hk_sol}
\end{align}
$H,K$在曲面保持定向的参数变换下不变。

\begin{theorem}
正则参数曲面的主曲率$\kappa_1,\kappa_2$是曲面上的连续函数；在任意非脐点的邻域内，$\kappa_1,\kappa_2$是连续可微函数。
\end{theorem}
\begin{proof}

曲面的一、二类基本量由参数方程二次偏导得到，故$H,K$均为连续可微函数。由\eqref{eq:principal_curvature_hk_sol}，主曲率连续。
非脐点满足$\kappa_1\neq\kappa_2$，即$H^2-K>0$，此时$\sqrt{H^2-K}$连续可微，故主曲率连续可微.

\end{proof}


\begin{enumerate}[(1)]
\item \textbf{非脐点情形}：$\kappa_1\neq\kappa_2$，将$\kappa_1,\kappa_2$依次代入\eqref{eq:principal_dir_linear_eq}，对应主方向满足
\begin{align}
\frac{\delta u}{\delta v}=-\frac{M-\kappa_1F}{L-\kappa_1E}=-\frac{N-\kappa_1G}{M-\kappa_1F},\label{eq:principal_dir_k1}\\
\frac{\delta u}{\delta v}=-\frac{M-\kappa_2F}{L-\kappa_2E}=-\frac{N-\kappa_2G}{M-\kappa_2F}.\label{eq:principal_dir_k2}
\end{align}
\item \textbf{脐点情形}：$\kappa_1=\kappa_2=\frac{L}{E}=\frac{M}{F}=\frac{N}{G}$，此时方程组系数矩阵为零矩阵，任意非零方向$(\delta u,\delta v)$均为主方向。
\end{enumerate}

由\eqref{eq:principal_dir_linear_eq}可得
\begin{align*}
\lambda=\frac{L\delta u+M\delta v}{E\delta u+F\delta v}=\frac{M\delta u+N\delta v}{F\delta u+G\delta v},
\end{align*}
因此主方向$(\mathrm{d}u,\mathrm{d}v)$满足行列式方程
\begin{align}
\begin{vmatrix}
L\mathrm{d}u+M\mathrm{d}v & E\mathrm{d}u+F\mathrm{d}v\\
M\mathrm{d}u+N\mathrm{d}v & F\mathrm{d}u+G\mathrm{d}v
\end{vmatrix}=0,\label{eq:principal_dir_det_eq}
\end{align}
展开为
\begin{align}
(LF-ME)(\mathrm{d}u)^2+(LG-NE)\mathrm{d}u\mathrm{d}v+(MG-NF)(\mathrm{d}v)^2=0.\label{eq:principal_dir_quad_eq}
\end{align}
该方程为曲面\textbf{曲率线}的微分方程，其解为主方向。改写为对称行列式形式：
\begin{align}
\begin{vmatrix}
(\mathrm{d}v)^2 & -\mathrm{d}u\mathrm{d}v & (\mathrm{d}u)^2\\
E & F & G\\
L & M & N
\end{vmatrix}=0.\label{eq:curvature_line_diff_eq}
\end{align}

\begin{theorem}
曲面$S:\boldsymbol{r}=\boldsymbol{r}(u,v)$上，参数曲线方向为彼此正交的主方向，当且仅当该点满足$F=M=0$。此时$u$-曲线方向主曲率为$\kappa_1=\frac{L}{E}$，$v$-曲线方向主曲率为$\kappa_2=\frac{N}{G}$。
\end{theorem}
\begin{proof}

必要性：$\boldsymbol{r}_u,\boldsymbol{r}_v$正交，故$F=0$。$(\mathrm{d}u,\mathrm{d}v)=(1,0)$为主方向，代入\eqref{eq:principal_dir_quad_eq}得$EM=0$，故$M=0$。
充分性：若$F=M=0$，则\eqref{eq:principal_dir_quad_eq}化为$(EN-GL)\mathrm{d}u\mathrm{d}v=0$，故$(\mathrm{d}u,\mathrm{d}v)=(1,0),(0,1)$均为主方向。代入主曲率公式得$\kappa_1=\frac{L}{E},\kappa_2=\frac{N}{G}$.

\end{proof}

\begin{corollary}
曲面的参数曲线网为正交曲率线网的充要条件是$F=M=0$。此时曲面的基本形式为
\begin{align}
\mathrm{I}=E(\mathrm{d}u)^2+G(\mathrm{d}v)^2,\quad \mathrm{II}=\kappa_1E(\mathrm{d}u)^2+\kappa_2G(\mathrm{d}v)^2.\label{eq:ortho_curvature_line_form}
\end{align}
\end{corollary}

\begin{theorem}
正则曲面的任意非脐点的邻域内，存在参数系$(u,v)$使得参数曲线构成正交曲率线网。
\end{theorem}
\begin{proof}

非脐点集为开集，其上主曲率连续可微，且存在两个正交的主方向场。由常微分方程理论，存在参数系使得参数曲线沿主方向，即构成正交曲率线网.

\end{proof}


Weingarten映射满足$W(\boldsymbol{r}_u)=-\boldsymbol{n}_u,\ W(\boldsymbol{r}_v)=-\boldsymbol{n}_v$，设
\begin{align}
\begin{pmatrix}
-\boldsymbol{n}_u\\
-\boldsymbol{n}_v
\end{pmatrix}=
\begin{pmatrix}
a_{11} & a_{12}\\
a_{21} & a_{22}
\end{pmatrix}
\begin{pmatrix}
\boldsymbol{r}_u\\
\boldsymbol{r}_v
\end{pmatrix},\label{eq:weingarten_matrix_def}
\end{align}
两边与$(\boldsymbol{r}_u,\boldsymbol{r}_v)$作内积，得
\begin{align*}
\begin{pmatrix}
L & M\\
M & N
\end{pmatrix}=
\begin{pmatrix}
a_{11} & a_{12}\\
a_{21} & a_{22}
\end{pmatrix}
\begin{pmatrix}
E & F\\
F & G
\end{pmatrix},
\end{align*}
因此
\begin{align}
\begin{pmatrix}
a_{11} & a_{12}\\
a_{21} & a_{22}
\end{pmatrix}=
\begin{pmatrix}
L & M\\
M & N
\end{pmatrix}
\begin{pmatrix}
E & F\\
F & G
\end{pmatrix}^{-1}.\label{eq:weingarten_matrix_formula}
\end{align}
二阶矩阵逆满足
\begin{align*}
\begin{pmatrix}
E & F\\
F & G
\end{pmatrix}^{-1}=\frac{1}{EG-F^2}
\begin{pmatrix}
G & -F\\
-F & E
\end{pmatrix},
\end{align*}
代入得
\begin{align}
\begin{pmatrix}
a_{11} & a_{12}\\
a_{21} & a_{22}
\end{pmatrix}=\frac{1}{EG-F^2}
\begin{pmatrix}
LG-MF & -LF+ME\\
MG-NF & -MF+NE
\end{pmatrix}.\label{eq:weingarten_matrix_expand}
\end{align}
线性变换的迹与行列式为不变量，故
\begin{align*}
2H=a_{11}+a_{22}=\frac{LG-2MF+NE}{EG-F^2},\quad K=a_{11}a_{22}-a_{12}a_{21}=\frac{LN-M^2}{EG-F^2}.
\end{align*}

由$\boldsymbol{n}_u\times\boldsymbol{n}_v=K\boldsymbol{r}_u\times\boldsymbol{r}_v$，取模长得
\begin{align}
|\boldsymbol{n}_u\times\boldsymbol{n}_v|=|K|\cdot|\boldsymbol{r}_u\times\boldsymbol{r}_v|.\label{eq:gauss_curvature_cross_prod}
\end{align}
曲面面积元素$\mathrm{d}\sigma=|\boldsymbol{r}_u\times\boldsymbol{r}_v|\mathrm{d}u\mathrm{d}v$，Gauss映射下的面积元素$\mathrm{d}\sigma_0=|\boldsymbol{n}_u\times\boldsymbol{n}_v|\mathrm{d}u\mathrm{d}v$，故$\mathrm{d}\sigma_0=|K|\mathrm{d}\sigma$。
设$D$为曲面围绕点$p$的邻域，$g(D)$为其Gauss映射的像，由积分中值定理，令$D\to p$取极限得
\begin{align}
|K(p)|=\lim\limits_{D\to p}\frac{A(g(D))}{A(D)}.\label{eq:gauss_curvature_geo_meaning}
\end{align}
即Gauss曲率的绝对值是Gauss映射下像区域与原区域的面积比在区域收缩至点$p$时的极限。

\begin{definition}
设$g:S\to\Sigma$为曲面$S$到单位球面$\Sigma$的Gauss映射，$\mathrm{I}_0$为$\Sigma$的第一基本形式，则$g^*\mathrm{I}_0$称为曲面$S$的\textbf{第三基本形式}，记为$\mathrm{III}$，即
\begin{align}
\mathrm{III}=g^*\mathrm{I}_0=\mathrm{d}\boldsymbol{n}\cdot\mathrm{d}\boldsymbol{n}=e(\mathrm{d}u)^2+2f\mathrm{d}u\mathrm{d}v+g(\mathrm{d}v)^2,\label{eq:third_form_def}
\end{align}
其中$e=\boldsymbol{n}_u\cdot\boldsymbol{n}_u,\ f=\boldsymbol{n}_u\cdot\boldsymbol{n}_v,\ g=\boldsymbol{n}_v\cdot\boldsymbol{n}_v$。
\end{definition}

\begin{theorem}
曲面的三个基本形式满足关系式
\begin{align}
\mathrm{III}-2H\mathrm{II}+K\mathrm{I}\equiv0.\label{eq:three_form_relation}
\end{align}
\end{theorem}
\begin{proof}

$\mathrm{I},\mathrm{II},\mathrm{III},H,K$均在定向参数变换下不变，故取点$p$附近正交曲率线参数系，满足$F=M=0$，且
\begin{align*}
\kappa_1=\frac{L}{E},\ \kappa_2=\frac{N}{G},\quad -\boldsymbol{n}_u=\kappa_1\boldsymbol{r}_u,\ -\boldsymbol{n}_v=\kappa_2\boldsymbol{r}_v.
\end{align*}
此时
\begin{align*}
e=\kappa_1^2E,\ f=0,\ g=\kappa_2^2G,
\end{align*}
三个基本形式为
\begin{align*}
\mathrm{I}=E(\mathrm{d}u)^2+G(\mathrm{d}v)^2,\quad \mathrm{II}=\kappa_1E(\mathrm{d}u)^2+\kappa_2G(\mathrm{d}v)^2,\quad \mathrm{III}=\kappa_1^2E(\mathrm{d}u)^2+\kappa_2^2G(\mathrm{d}v)^2.
\end{align*}
代入验证：
\begin{align*}
2H\mathrm{II}-K\mathrm{I}=(\kappa_1+\kappa_2)\left(\kappa_1E(\mathrm{d}u)^2+\kappa_2G(\mathrm{d}v)^2\right)-\kappa_1\kappa_2\left(E(\mathrm{d}u)^2+G(\mathrm{d}v)^2\right)=\kappa_1^2E(\mathrm{d}u)^2+\kappa_2^2G(\mathrm{d}v)^2=\mathrm{III},
\end{align*}
即$\mathrm{III}-2H\mathrm{II}+K\mathrm{I}\equiv0$.

\end{proof}

\begin{example}
求螺旋面$\bm{r}=(u\cos v,u\sin v,u+v)$的Gauss曲率和平均曲率。
\end{example}

\begin{example}
设$S$是平面曲线$C$:$\bm{r}=(g(s),0,f(s))$围绕$z$轴旋转一周所得的曲面，其中$s$是曲线$C$的弧长参数。求曲面$S$的Gauss曲率$K$。
\end{example}

\begin{example}
求下列曲面的Gauss曲率$K$和平均曲率$H$:
\begin{enumerate}[(1)]
\item $\bm{r}=(u\cos v,u\sin v,ku^{2})$;
\item 曲线$z=k\sqrt{x}$,$y=0$绕$z$轴旋转所生成的曲面;
\item 曲线(曳物线)
\begin{align*}
\bm{r}=\left(k\sin v,0,k\left(\ln\tan\frac{v}{2}+\cos v\right)\right)
\end{align*}
绕$z$轴旋转所生成的曲面(伪球面);
\item Scherk曲面$z=\frac{1}{k}(\ln(\cos kx)-\ln(\cos ky))$.
\end{enumerate}
\end{example}

\begin{example}
求双曲抛物面$\bm{r}=(a(u+v),b(u-v),2uv)$的Gauss曲率$K$，平均曲率$H$，主曲率$\kappa_1,\kappa_2$和它们所对应的主方向。
\end{example}

\begin{example}
设在曲线$C$:$\bm{r}=\bm{r}(s)$的所有法线上截取长度为$\lambda$的一段，它们的端点的轨迹构成一个曲面$S$，称为围绕曲线$C$的管状曲面，其方程是
\begin{align*}
\bm{r}(s,\theta)=\bm{r}(s)+\lambda(\cos\theta\bm{\beta}(s)+\sin\theta\bm{\gamma}(s)),
\end{align*}
其中$s$是曲线的弧长参数，$\bm{\beta}(s),\bm{\gamma}(s)$分别是曲线$C$的主法向量和次法向量。求该曲面的Gauss曲率$K$，平均曲率$H$和主曲率$\kappa_1,\kappa_2$。
\end{example}

\begin{example}
在曲面$S$:$\bm{r}=\bm{r}(u,v)$上每一点沿法线截取长度为$\lambda$(适当小的正数)的一段，它们的端点的轨迹构成一个曲面$\tilde{S}$，称为原曲面$S$的平行曲面，其方程是
\begin{align*}
\tilde{\bm{r}}(u,v)=\bm{r}(u,v)+\lambda\bm{n}(u,v).
\end{align*}
从点$\bm{r}(u,v)$到点$\tilde{\bm{r}}(u,v)$的对应记为$\sigma$。
\begin{enumerate}[(1)]
\item 曲面$S$和曲面$\tilde{S}$在对应点的切平面互相平行。

\item 对应$\sigma$把曲面$S$上的曲率线映为曲面$\tilde{S}$上的曲率线。

\item 曲面$S$和曲面$\tilde{S}$在对应点的Gauss曲率和平均曲率有下列关系:
\begin{align*}
\tilde{K}=\frac{K}{1-2\lambda H+\lambda^{2}K},\quad \tilde{H}=\frac{H-\lambda K}{1-2\lambda H+\lambda^{2}K}.
\end{align*}
\end{enumerate}
\end{example}

\begin{example}
求曲面$z=f(x,y)$上的脐点所满足的(充分必要)条件。
\end{example}

\begin{example}
求下列曲面上的脐点:
\begin{enumerate}[(1)]
\item $2z=\frac{x^{2}}{\alpha^{2}}+\frac{y^{2}}{\beta^{2}}$,其中$\alpha>\beta>0$;
\item $\frac{x^{2}}{\alpha^{2}}+\frac{y^{2}}{\beta^{2}}+\frac{z^{2}}{\gamma^{2}}=1$,其中$\alpha>\beta>\gamma>0$.
\end{enumerate}
\end{example}

\begin{example}
证明下列恒等式:
\begin{enumerate}[(1)]
\item \begin{align*}
\begin{pmatrix} L&M\\ M&N \end{pmatrix}\cdot \begin{pmatrix} E&F\\ F&G \end{pmatrix}^{-1}
=\frac{1}{\sqrt{EG-F^{2}}}\begin{pmatrix} -\bm{n}_{u}\cdot(\bm{r}_{v}\times\bm{n})&\bm{n}_{u}\cdot(\bm{r}_{u}\times\bm{n})\\ -\bm{n}_{v}\cdot(\bm{r}_{v}\times\bm{n})&\bm{n}_{v}\cdot(\bm{r}_{u}\times\bm{n}) \end{pmatrix};
\end{align*}

\item \begin{align*}
\begin{pmatrix} L&M\\ M&N \end{pmatrix}\cdot \begin{pmatrix} E&F\\ F&G \end{pmatrix}^{-1}\cdot \begin{pmatrix} L&M\\ M&N \end{pmatrix}=\begin{pmatrix} e&f\\ f&g \end{pmatrix}.
\end{align*}
\end{enumerate}
\end{example}



\end{document}